\begin{conf-abstract}
{Seasonality of Excess Air Signatures Under Different Flow Conditions}
{Christoph West, Nicholas Thiros, Werner Aeschbach, Payton Gardner}
{Institute of Environmental Physics Heidelberg}
{'- hydrograph of a high altitude fractured shale bedrock mountain aquifer in the East River watershed, Colorado, US driven by seasonal snowmelt patterns - results in two flow periods: recharge (peak snowmelt, water table rise), baseflow (no snowmelt/recharge events, water table lower at a base level) - bulk noble gas measurements during recharge period show signature of excess air, while measurements during the baseflow period show a degassed signature - subsurface production of CO2, N2, and CH4 have been observed previously resulting from bedrock weathering of the shale - Possible explanation: - Higher water table (during the recharge period) corresponds to higher hydrostatic pressure, which can dilute more gas into the groundwater, thus showing excess air in the noble gases - Lower water table (during baseflow) corresponds to lower hydrostatic pressure, thus oversaturated gas exsolves continuously and leads to the degassed signature}
\end{conf-abstract}
